{\scriptsize
\begin{flushleft}
\end{flushleft}

        \begin{tabularx}{\linewidth}{|p{0.35\textwidth}|X|}
            \hline 
            \textbf{Notions et contenus} & \textbf{Capacités exigibles} \\
            \hline 
            \multicolumn{2}{|l|}{\textbf{5.3.2. Exemples de champs magnétostatiques }} \\
            \hline 
            Solénoïde long sans effets de bords. &  Solénoïde long sans effets de bords. \\ 
            \hline
            Inductance propre. Densité volumique d'énergie magnétique. & Établir les expressions de l'inductance propre et de l'énergie d'une bobine modélisée par un solénoïde long. Associer l'énergie d'une bobine à une densité volumique d'énergie magnétique. \\
            \multicolumn{2}{|l|}{\textbf{5.3.3. Dipôles magnétostatiques}} \\
            \hline 
            Moment magnétique d'une boucle de courant plane.  & 
            Relier le moment magnétique d'un atome d'hydrogène à son moment cinétique. \\
            \hline 
            Rapport gyromagnétique de l'électron. Magnéton de Bohr.  & Construire en ordre de grandeur le magnéton de Bohr par analyse dimensionnelle.
            Évaluer l'ordre de grandeur maximal du moment magnétique volumique d'un aimant permanent.  \\
            \hline 
            Actions subies par un dipôle magnétique placé dans un champ magnétostatique d'origine extérieure : résultante et moment. & Utiliser les expressions fournies de la résultante et du moment des actions subies par un dipôle magnétique placé dans un champ magnétostatique d'origine extérieure Décrire l'expérience de Stern et Gerlach et expliquer ses enjeux. \\
            \hline
            Énergie potentielle d'un dipôle magnétique rigide placé dans un champ magnétostatique d'origine extérieure.  & Utiliser l'expression fournie de l'énergie potentielle d'un dipôle rigide dans un champ magnétostatique d'origine extérieure.
            Prévoir qualitativement l'évolution d'un dipôle rigide dans un champ magnétostatique d'origine extérieure. \\
        \hline 
        \end{tabularx}
\vspace{0.5cm}
}
